\documentclass[11pt]{article}
\usepackage[a4paper,margin=1in]{geometry}
\usepackage{amsmath,amssymb,amsthm}
\usepackage{booktabs}
\usepackage{graphicx}
\usepackage{xcolor}
\usepackage{hyperref}
\hypersetup{colorlinks=true,linkcolor=blue,urlcolor=blue,citecolor=blue}
\newcommand{\rt}[1]{\texttt{#1}}
\newcommand{\fn}[1]{\texttt{#1}}
\newcommand{\ifgraphic}[2]{\IfFileExists{#1}{\includegraphics[width=#2]{#1}}%
  {\fbox{\parbox{#2}{\centering\small figure \texttt{\detokenize{#1}} pending}}}}
\newcommand{\iftable}[1]{\IfFileExists{#1}{\input{#1}}%
  {\fbox{\parbox{0.9\textwidth}{\centering\small table
  \texttt{\detokenize{#1}} pending}}}}

\theoremstyle{plain}
\newtheorem{proposition}{Proposition}
\newtheorem{theorem}{Theorem}

\title{\textbf{Report R21 --- strong or weak? Grading the fragmentation of
108/201 and 156/198}\\
\large An addendum to R18 --- Krylov sector analysis of quantum cellular
automata}
\author{Tier 1 analysis}
\date{2026-08-14}

\begin{document}
\maketitle

\begin{abstract}
R9 and R19 report $D_{\max}/2^N \to 0$ for all four rules R18 is about. That is
necessary for fragmentation but it is not the criterion in use: the literature
grades fragmentation by the largest Krylov sector relative to the
\emph{symmetry sector} that contains it, $D_{\max}(s)/\dim s$, and dividing by
$2^N$ instead is the same test only when the symmetry is trivial --- which here
it is not. We run the proper test. \textbf{All four rules are strongly
fragmented, and the conclusion is not close.} At $N=20$ the ratio in the
dominant symmetry sector is $5.9\times10^{-3}$ ($W_{108}$),
$8.8\times10^{-3}$ ($W_{201}$) and $3.0\times10^{-3}$ ($W_{156}/W_{198}$),
decaying geometrically at rates $0.841$, $0.834$ and $0.697$ against the
predicted $\varphi/2 = 0.809$ and $4^{1/5}/2 = 0.660$. We do not take R18's
charges on trust: a detector built here finds the \emph{complete} space of
conserved charges $\sum_i q(x_i,\ldots,x_{i+r-1})$ by exact rational
elimination, at ranges $r=1,2,3,4$ and on a two-site unit cell, and
independently reproduces R18's pair at $r=2$ --- for $W_{108}$ it returns
literally the sublattice-resolved counts of the wall word \rt{00}, with no
notion of a wall anywhere in its input. Longer range finds \emph{more} charges
(ten for $W_{108}$ at $r=4$, which R18 did not report) and the verdict does not
move. Underneath the numerics there is a theorem: finitely many finite-range
extensive charges have at most $O(N^m)$ joint level sets, so the symmetry can
only ever remove a \emph{polynomial} factor, and any rule whose $D_{\max}$
grows like $\lambda^N$ with $\lambda<2$ is strongly fragmented for that reason
alone. Finally, the diagnostic is shown not to be vacuous on this project's own
rules: $W_{60}/W_{102}$ hold exactly half the Hilbert space in one sector and
grade \textbf{weak}, while $W_{150}$ and $W_{105}$ turn out to be
\textbf{not fragmented at all} --- their Krylov sectors are exactly the level
sets of a single conserved charge, the total domain-wall number $D$ for
$W_{150}$ and the staggered one $S$ for $W_{105}$. Everything is repeated at
\rt{pbc} to $N=18$ with the same verdict, and there the detected charge level
sets reproduce R18's Tier-2 counts ($13,18,25$ and $12,17,23$) exactly.
\end{abstract}

\tableofcontents

\section{The question, stated carefully}

\subsection{Two ratios, one of which is the criterion}

Write $\mathcal{K}(x)$ for the Krylov sector of a computational basis state $x$
and $s(x)$ for the symmetry sector containing it. The grading used in the
literature (Sala \emph{et al.}, PRX \textbf{10}, 011047 (2020); Khemani,
Hermele and Nandkishore, PRB \textbf{101}, 174204 (2020); Moudgalya, Bernevig
and Regnault, Rep.\ Prog.\ Phys.\ \textbf{85}, 086501 (2022)) is
\[
\textbf{strong:}\quad \frac{D_{\max}(s)}{\dim s}\;\longrightarrow\;0,
\qquad\qquad
\textbf{weak:}\quad \frac{D_{\max}(s)}{\dim s}\;\longrightarrow\;\text{const}>0 ,
\]
in the thermodynamic limit, with $s$ the dominant (largest) symmetry sector.
Weak fragmentation still has one thermal Krylov sector that fills its symmetry
sector, with the frozen and small sectors a vanishing remainder; strong
fragmentation has no such sector, and a typical initial state explores an
asymptotically vanishing fraction of what its conservation laws allow.

R9's $D_{\max}/2^N$ is the same quantity only if $\dim s = 2^N$, i.e.\ only if
there is no symmetry to speak of. For these four rules there is: R18 certifies
two independent local charges apiece. So the crude ratio understates
$\dim s$ and therefore \emph{overstates} how strongly fragmented the model is.
The question is whether it overstates it enough to matter.

\subsection{The trap: ``all conserved quantities'' makes the question empty}

Taken literally, the grading is not well posed. R18 proves that the
\emph{wall-occurrence set} is conserved and is a \emph{complete invariant}: two
states share a Krylov sector if and only if the rule's minimal wall word occurs
at the same positions. If one is allowed to call that a symmetry, the symmetry
sectors \emph{are} the Krylov sectors, the ratio is $1$ by construction, and
every fragmented model in the world is ``weak''.

The grading is therefore always relative to a \emph{declared} symmetry algebra
--- the local charges one would write down for the model --- held at fixed
finite range while $N\to\infty$. The honest procedure is not to declare as few
charges as possible but as \emph{many} as exist at each range, since a larger
symmetry makes $\dim s$ smaller and the ratio larger, and hence makes a
``strong'' verdict harder to earn. That is what this report does, and
Section~\ref{sec:detector} is how.

\section{Finding the charges rather than assuming them}
\label{sec:detector}

\subsection{The ansatz and the elimination}

\fn{scaling.hsf\_strength} detects every charge of the form
\[
Q(x) \;=\; \sum_i q\bigl(i \bmod p;\; s_i, s_{i+1}, \ldots, s_{i+r-1}\bigr),
\]
with $s$ the padded string the dynamics actually sees (\rt{obc0} pads a frozen
$0$ outside each end, and those zeros are charge material for the same reason
R18 found them to be wall material). The method is: label every state with its
Krylov sector via the flip reduction (\fn{graph.flip\_graph}), form the
count-vector difference between each state and a representative of its sector,
and take the exact rational null space of the resulting row space. Every basis
vector is then re-evaluated and \emph{asserted} to be constant on every sector,
so nothing rests on floating point.

Two conventions are not free choices.

\paragraph{The unit cell must be two sites.} A brick-wall cycle is invariant
under two-site translations only, so a conserved density may live on one
sublattice. With $p=1$ the detector finds exactly \emph{one} charge for each of
the four rules; with $p=2$ it finds \emph{two}. The second is the staggered one,
and $p=2$ is R18's ``parity split'' --- the reason its Tier-2 detector needed
one too.

\paragraph{Detection must be parity-resolved in $N$.} On an open chain the two
sublattices swap roles with $N \bmod 2$, so a basis found at $N=12$ need not be
conserved at $N=13$. It happens to survive at $r=2$, which is how an earlier
pass of this module got away with transferring one basis to every $N$; at
$r=4$ every rule in play has at least one basis vector that breaks. The module
detects separately for even and odd $N$, and \fn{strength} reports
\rt{krylov\_inside\_sym} at every measured point --- the flag that caught it,
and which is asserted throughout.

\subsection{It reproduces R18, from a completely different direction}

\begin{proposition}
At $r=2$, $p=2$ the detector finds exactly two independent charges for each of
$W_{108}, W_{201}, W_{156}, W_{198}$, and their joint level sets coincide
exactly with R18's $(|W_{\rm even}|, |W_{\rm odd}|)$ --- the sublattice-resolved
occurrence counts of the rule's minimal wall word. Verified for $N=10,11,12$.
\end{proposition}

\noindent For $W_{108}$ the agreement is more than a coincidence of partitions:
the returned integer basis is, verbatim,
\[
q_1 = \bigl[\text{even positions: } \rt{00}\mapsto 1\bigr], \qquad
q_2 = \bigl[\text{odd positions: } \rt{00}\mapsto 1\bigr],
\]
i.e.\ $|W_{\rm even}|$ and $|W_{\rm odd}|$ for the word \rt{00}, which is
exactly the word R18 reports. The detector was given the sector partition and
a pattern-counting ansatz and nothing else; it recovered the wall word. For the
other three rules the basis is a different (equally valid) choice of coordinates
on the same two-dimensional space, and the level sets agree exactly.

This is the cross-check R18 could not perform on itself: its walls and its
charges both came from Tier-2 modules, whereas this detector is Tier-1 code
with no wall concept in it.

\subsection{What R18 did not report: longer range finds more}

R18's ``exactly two independent local charges'' is a statement about range 2.
It does not survive to range 3 or 4:

\begin{center}
\small
\begin{tabular}{lcccc}
\toprule
& $r=1$ & $r=2$ & $r=3$ & $r=4$ \\
\midrule
$W_{108}$ & 0 & 2 & 6 & 10 \\
$W_{201}$ & 0 & 2 & 4 & 8 \\
$W_{156}$, $W_{198}$ & 0 & 2 & 3 & 7 \\
\midrule
$W_{60}$, $W_{102}$ & 0 & 0 & 1 & 2 \\
$W_{150}$, $W_{105}$ & 0 & 1 & 1 & 1 \\
\bottomrule
\end{tabular}
\end{center}

\noindent The $r=1$ column is R18's other observation, now checked
independently: there is no single-site charge for any of them, so magnetisation
is not conserved --- immediate once one notices the gate flips a site outright.

The growth of the other columns is what makes this section necessary rather than
decorative. Ten charges could in principle carve the Hilbert space finely enough
to change the verdict, and the only way to know is to measure at each range
separately. Section~\ref{sec:results} does, and Section~\ref{sec:theorem}
explains why the answer was never in doubt.

\section{The measurements}
\label{sec:results}

\begin{center}
\footnotesize
\iftable{tab_r21_strength_obc0.tex}
\end{center}

\begin{figure}[htbp]
\centering
\ifgraphic{../../figures/fig_hsf_strength_obc0.pdf}{\textwidth}
\caption{\textbf{The crude ratio and the criterion give the same verdict, for a
reason.} (a) $D_{\max}/2^N$, as R9 reports it. (b) The criterion,
$D_{\max}(s^*)/\dim s^*$ in the dominant symmetry sector at range 2; the
zig-zag on the four fragmented curves is the even/odd split of the chain.
(c) The average over computational basis states of $D_K(x)/\dim s(x)$ --- the
fraction of its own symmetry sector a random product state can reach --- which
is the same statement without the parity artefact. Dashed guides are
$(\varphi/2)^N$ and $(4^{1/5}/2)^N$. $W_{60}$ and $W_{102}$ sit flat at $1/2$
(weak); $W_{150}$ and $W_{105}$ sit flat at $1$ (not fragmented at all).}
\label{fig:strength}
\end{figure}

\noindent The numbers behind panel (b), at range 2:

\begin{center}
\small
\iftable{tab_r21_ratios_obc0.tex}
\end{center}

\noindent Three things to read off.

\paragraph{The symmetry matters, and does not save anything.} At $N=20$,
$W_{156}$'s crude ratio is $3.1\times10^{-4}$ while the proper one is
$3.0\times10^{-3}$ --- a factor of ten. So dividing by $2^N$ genuinely was the
wrong denominator. It is also a factor of ten and not a factor of $2^{cN}$,
which is why the verdict is unchanged. For $W_{108}$ the two ratios are
$6.5\times10^{-3}$ and $5.9\times10^{-3}$, nearly equal.

\paragraph{The dominant symmetry sector holds $\Theta(2^N/N)$ states.} The
quantity $\dim(s^*)\,N/2^N$ sits at $1.56, 1.79, 1.83, 1.97$ for $W_{108}$ over
$N=8,12,16,20$ and similarly for the others --- flat to within a factor, which
is the numerical content of ``the symmetry only removes a polynomial factor''.
The number of symmetry sectors itself runs $15, 28, 45, 66$ at $N=8,12,16,20$,
i.e.\ $\binom{N/2+2}{2}$, visibly quadratic.

\paragraph{The decay rate is the predicted one.} With no charges at all
($r=1$) the measured rate is $0.8090$ for $W_{108}$, against $\varphi/2 =
0.809017$, and $0.6581$ for $W_{156}$ against $4^{1/5}/2 = 0.659754$ --- six and
three digits respectively, with the exponent derived (R18's wall recurrence,
R19's closed forms, R20's OEIS entries) rather than fitted. Switching the
charges on lifts the rate slightly, to $0.841$ and $0.697$, which is exactly the
$O(N^2)$ prefactor showing up as a slower apparent decay over a finite window.

\subsection{Every range, separately}

\begin{center}
\small
\iftable{tab_r21_ranges_obc0.tex}
\end{center}

\noindent All four rules grade \textbf{strong} at $r=1,2,3$ and $4$. One caveat
belongs with the $r=4$ rows and is stated rather than buried: with $m=10$
charges the bound below carries an $N^{10}$ prefactor, and $N\le20$ is nowhere
near the asymptotic regime for such a prefactor. The $r=4$ ratios are therefore
noisy and occasionally sit at $1$ at small $N$. They are reported because
suppressing them would be worse; the argument that settles $r=4$ is the theorem,
not the fit.

\section{Why the answer could not have been otherwise}
\label{sec:theorem}

\begin{theorem}
Fix a range $r$ and unit cell $p$, and let $Q_1,\ldots,Q_m$ be conserved charges
of the above form, $m \le p\,2^r$. Then the number of joint level sets is
$O(N^m)$, and consequently
\[
\frac{D_{\max}(s^*)}{\dim s^*} \;=\; O\!\left(N^{m}\,\frac{D_{\max}}{2^N}\right).
\]
If $D_{\max} = \Theta(\lambda^N)$ with $\lambda<2$, the ratio vanishes
exponentially, at rate $\lambda/2$ up to the polynomial factor, and the
fragmentation is strong.
\end{theorem}

\begin{proof}
Each $Q_j$ is a sum of at most $L = N+2$ terms drawn from the finite set of
values of $q_j$, so $Q_j$ takes $O(N)$ distinct values and the $m$ of them have
$O(N^m)$ joint level sets. The largest of those level sets therefore holds at
least $2^N/O(N^m)$ states, and $D_{\max}(s^*) \le D_{\max}$.
\end{proof}

\noindent Three consequences worth stating plainly.

\begin{enumerate}
\item \textbf{The user's instinct was right, and the reason is this theorem.}
``$D_{\max}/2^N \to 0$, so it should be strong'' is not a sloppy shortcut when
the decay is exponential: a finite-range symmetry cannot supply more than a
polynomial denominator, so the two criteria agree. What the theorem also shows
is where the shortcut fails --- when the decay is only polynomial, the symmetry
can absorb all of it, and the two criteria come apart. Both cases occur in this
project (Section~\ref{sec:controls}).

\item \textbf{Range cannot rescue the symmetry story, at any fixed $r$.} Ten
charges give $N^{10}$, which is enormous at $N=20$ and irrelevant at $N=200$.
The only way a symmetry could compete is if $m$ grew with $N$ --- i.e.\ if the
declared symmetry were not of fixed finite range --- and that is precisely the
degenerate case of Section~1.2, in which the wall set itself is the symmetry.

\item \textbf{It sharpens what R18 proved.} R18 showed the charges are too
\emph{coarse} to explain the sectors: $O(N^2)$ level sets against $\rho^{2N}$ or
$\varphi^N$ Krylov sectors, the ratio growing exponentially. That is a statement
about counting sectors. The theorem here is the corresponding statement about
their \emph{sizes}, and it is what turns ``the symmetry is not the explanation''
into ``the fragmentation is strong''.
\end{enumerate}

\section{The diagnostic is not vacuous: three grades among eight rules}
\label{sec:controls}

A test that returns ``strong'' for everything would say nothing. Running the
same pipeline over the polynomial corner of the unitary sweep produces three
distinct outcomes, all within this project's own rules.

\paragraph{$W_{60}$, $W_{102}$ --- weak.} $D_{\max}=2^{N-1}$: half the Hilbert
space sits in one Krylov sector, at every $N$. There is no charge at all below
range 3, so the symmetry sector is the whole space and
$D_{\max}(s^*)/\dim s^* = 1/2$ exactly, forever. The average explored fraction
is exactly $1/3$. This is the textbook weak picture --- one thermal sector
filling its symmetry sector, with $N$ small sectors beside it --- realised by a
rule of this project rather than a model imported for the comparison.

\paragraph{$W_{150}$, $W_{105}$ --- not fragmented at all.} Both have exactly
\emph{one} conserved range-2 charge, and its level sets \emph{are} their Krylov
sectors: $n_{\rm sym} = n_{\rm Krylov} = 11$ at $N=20$, ratio identically $1$,
explored fraction identically $1$, at every $N$ and every range. Calling this
``weak'' would be wrong; weak fragmentation still has small sectors beside the
thermal one, and these rules have none. The charges are named:

\begin{proposition}
$W_{150}$ conserves the total domain-wall number $D = \sum_i b_i$ and not the
staggered $S$; $W_{105}$ conserves $S = \sum_i (-1)^i b_i$ and not $D$. In each
case the level sets of that single charge are exactly the Krylov sectors.
\end{proposition}

\noindent This retro-explains two things. It explains why $W_{150}$'s sector
count is the trivial $\lfloor (N+1)/2\rfloor + 1$: the \rt{obc0} padding forces
an even number of domain walls among the $N+1$ bonds, so $D$ takes exactly that
many values --- connecting to R8's domain-wall count for rule 150, and matching
R20's identification of the count as \rt{A008619}$(N+1)$. And it explains R20's identifications of their
$D_{\max}$: $W_{105}$'s largest sector is the central binomial coefficient
$\binom{N+1}{\lfloor (N+1)/2\rfloor}$ because that is the largest level set of a
balanced counting statistic, and $W_{150}$'s is \rt{A214282} for the same
reason. Those closed forms are the symmetry showing through, not accidents.

\paragraph{$W_{108}$, $W_{201}$, $W_{156}$, $W_{198}$ --- strong.} The four
rules of R18, at every range tested, by both the finite-$N$ measurement and the
theorem.

\medskip
\noindent The interesting structural fact is that the split is exactly the split
in $\lambda$: the strong rules are the ones with $\lambda = \varphi$ or
$4^{1/5}$, both $<2$; the weak and unfragmented ones all have
$\lambda = 2$, differing only in the sub-exponential correction
($2^{N-1}$ against $2^N/\sqrt{N}$). By the theorem that is not a coincidence but
the whole content of the classification.

\section{The same question on the ring}

R18 tested both boundary conditions, so this addendum does too. The pipeline is
exhaustive enumeration to $N=18$ --- the same kind of computation R18 already ran
at \rt{pbc}, not the production sweep.

\begin{center}
\footnotesize
\iftable{tab_r21_strength_pbc.tex}
\end{center}

\noindent The verdict is the same: \textbf{strong for all four rules at every
range}, with $D_{\max}(s^*)/\dim s^*$ at $7.6\times10^{-3}$ ($W_{108}$,
$W_{201}$) and $4.7\times10^{-3}$ ($W_{156}$, $W_{198}$) at $N=18$. Three
details are specific to the ring.

\paragraph{A third confirmation of R18, on the nose.} At range 2 the detected
charges have $13, 18, 25$ joint level sets for $W_{108}$ at $N=8,10,12$ and
$12, 17, 23$ for $W_{156}$ --- exactly the numbers R18 quotes for Tier-2's
\emph{full certified charge set} at \rt{pbc}. That was R18's own check that
nothing conserved had been left out; it now reproduces from a third,
independent detector.

\paragraph{$W_{108}$ and $W_{201}$ merge.} On the ring their series coincide
term for term (same charge counts, same sector counts, same ratios), because the
spin flip that relates them is a \rt{pbc} symmetry and not an \rt{obc0} one ---
consistent with R19's finding that the \rt{obc0} vacuum breaks it.

\paragraph{The controls move, and in the expected direction.}
$W_{60}$ and $W_{102}$ lose their charges entirely (none at any range up to 4)
and collapse to \emph{two} Krylov sectors of $2^{N-1}$ each, so the ratio is
$1$ rather than $1/2$; the periodic chain removes the boundary that split the
smaller sectors off. $W_{105}$ stays unfragmented ($9$ sectors, $9$ level sets
at $N=18$), while $W_{150}$ acquires a small excess ($12$ Krylov sectors inside
$10$ level sets) and grades weak rather than unfragmented. None of this touches
the four rules of R18.

\section{Answering the question as asked}

\begin{enumerate}
\item \textbf{Does $D_{\max}/2^N\to0$ imply strong HSF?} Not in general --- but
yes whenever the decay is exponential and the declared symmetry is of fixed
finite range, which covers all four rules here. The two criteria come apart
only in the polynomial case, and this project contains examples of that too.

\item \textbf{Do 108, 201, 156, 198 exhibit weak or strong HSF, measured against
the R18 charges?} \textbf{Strong}, at $N=20$, at ranges 1 through 4, with the
ratio at $5.9\times10^{-3}$, $8.8\times10^{-3}$, $3.0\times10^{-3}$ and
$3.0\times10^{-3}$ and decaying geometrically at the predicted rates.

\item \textbf{Were $S$, $D$, $n_{01}$, $n_{10}$ the right charges to divide
by?} They are a spanning set but not an independent one: for each rule the
independent pair at range 2 is the sublattice-resolved wall counts, of which
$D$ or $S$ is a function (R18, confirmed here by an independent detector).
Dividing by fewer would have given a larger $\dim s$ and an even smaller ratio,
so the choice made here is the conservative one. Dividing by \emph{more} --- the
range-3 and range-4 charges R18 did not report --- also leaves the verdict
unchanged.

\item \textbf{Is there anything the symmetry could have done?} No, and the
theorem says why: $m$ finite-range charges have $O(N^m)$ level sets, so no
choice of local symmetry can supply an exponential denominator. The only way to
close the gap is with the wall set itself, which is not a local charge and is
the entire point of R18.
\end{enumerate}

\section{Reproduce}

\begin{verbatim}
python -m qca_fragmentation.scaling.hsf_strength \
       --bc obc0 --rebuild --n-max 20
python -m qca_fragmentation.scaling.hsf_strength \
       --bc pbc  --rebuild --n-max 18
python -m pytest qca_fragmentation/tests/test_hsf_strength.py -q
\end{verbatim}

\noindent Everything here comes from \fn{analytics/hsf\_strength\_\{bc\}.json}.
The Krylov labels use the flip reduction of \fn{graph.flip\_graph}, valid
because all eight rules are unitary; the charge detector, the ratios and the
figure are \fn{scaling/hsf\_strength.py}, which depends on R18's
\fn{scaling/wall\_charges.py} only for the cross-check in \S2.2 and for naming
$D$ and $S$ in \S5.

\end{document}
